% ===================================================================
%  TABEL Nr. 3 — Kinderjare en jeug.
%  Meme regime que les tableaux 1 et 2.
%
%  « eerste tafereel » N'A RIEN DEMANDE A html.py : les deux mots
%  etaient deja dans le groupe commun, poses par le neerlandais. C'est
%  le seul endroit de cette colonne ou la parente avec texto/nl sert
%  au lieu de menacer.
%
%  LE DIMINUTIF EST PARTOUT DANS CE TABLEAU — wolkies, swaeltjies,
%  kuikentjies, dogtertjie — et il s'ecrit -tjie, -jie, -kie, -pie ou
%  -ie, jamais -tje. La regle de kolonoj.py le releve ; ce qu'elle ne
%  peut pas relever, c'est un diminutif OUBLIE la ou l'ido en met un.
%
%  ET LA NEGATION SE FERME : « nie ... nie ». Ce tableau n'en a que
%  deux, mais elles sont du meme piege que celles du tableau 2.
% ===================================================================

%%K t03-apar-1 apar
\VUsaut{3.0mm}
\VUtitre{15.0pt}{60}{TABEL Nr. 3}
\VUsaut{1.6mm}
\VUtitre{11.4pt}{0}{Kinderjare en jeug.}
\VUsaut{3.4mm}

%%K t03-apar-2 apar
(Tabelle 3 en 4 is so saamgestel dat hulle in vier \VUgras{gesinstonele} die wesenlike
begrippe oor die \VUgras{weer}, die \VUgras{ouderdom}, die \VUgras{gesin}, die
\VUgras{kleredrag} en die \VUgras{stand} aanbied.)
\VUsaut{4.0mm}

%%K t03-tit-1 sub
\VUcentre{12.0pt}{0}{\textit{Eerste tafereel.}}
\VUsaut{3.0mm}

%%K t03-tit-2 sub
\VUpk{10.2pt}{40}{Doopplegtigheid in 'n dorp.}
\VUsaut{2.4mm}

%%K t03-tit-3 sub
\VUpk{10.2pt}{40}{Die lente.}
\VUsaut{3.0mm}

%%K t03-c1-01-1 p
1. --- Ons is in die \VUgras{lente}, in die maande \VUgras{Maart}, \VUgras{April} of
\VUgras{Mei}, op die dae van die \VUgras{week}: \VUgras{Maandag}, \VUgras{Dinsdag} of
\VUgras{Woensdag}. Dit is \VUgras{tienuur} in die oggend.

%%K t03-c1-02-1 p
2. --- Die \VUgras{weer} is mooi, die \VUgras{lug} suiwer. Die son skyn. 'n Paar
\VUgras{wolkies} \textsuperscript{(2)} styg op in die blou \VUgras{hemel}
\textsuperscript{(1)}.

%%K t03-c1-03-1 p
3. --- Die \VUgras{bome} het nog skaars \VUgras{blare}. Die skraal \VUgras{populiere}
\textsuperscript{(3)} en die hoë \VUgras{plataan} \textsuperscript{(17)} het tot nou toe
net knoppe.

%%K t03-c1-04-1 p
4. --- Die \VUgras{swaeltjies} \textsuperscript{(4)} kom terug en vlieg met vrolike
getjirp om die \VUgras{spits} \textsuperscript{(7)} van die \VUgras{kloktoring}
\textsuperscript{(5)} wat die \VUgras{kerk} \textsuperscript{(6)} van die dorp bekroon.

%%K t03-tit-4 sub
\VUsaut{3.4mm}
\VUpk{10.2pt}{40}{Die kerk.}
\VUsaut{3.0mm}

%%K t03-c1-05-1 p
5. --- Baie eenvoudig is daardie kerk met sy groot \VUgras{portaal}
\textsuperscript{(13)} en dik \VUgras{horlosie} \textsuperscript{(8)}. Die klein
\VUgras{wyser} \textsuperscript{(9)} is op die syfer X, dit dui die \VUgras{ure} aan.
Die \VUgras{groot een} \textsuperscript{(10)} is op XII (middel); dit dui die
\VUgras{minute} aan.

%%K t03-c1-06-1 p
6. --- In die kloktoring lui die \VUgras{klokke} \textsuperscript{(11)} vrolik. Bo op
die dak staan 'n klein klip \VUgras{kruis} \textsuperscript{(12)}.

%%K t03-c1-07-1 p
7. --- Die kerk het nie net 'n groot \VUgras{portaal} \textsuperscript{(13)} nie, dit
het ook 'n sydeurtjie. Deur hierdie deurtjie kom meneer die \VUgras{pastoor}
\textsuperscript{(15)}, geklee in sy \VUgras{toga}, uit die \VUgras{sakristie}
\textsuperscript{(14)} en gaan na die \VUgras{pastorie} \textsuperscript{(19)} toe.

%%K t03-c1-08-1 p
8. --- Naby hom, hier is die \VUgras{melkvrou} \textsuperscript{(24)} met haar
\VUgras{donkie} \textsuperscript{(23)}. Sy kom terug van die stad, waarheen sy goeie melk
vir die kinders gebring het.

%%K t03-tit-5 sub
\VUsaut{4.0mm}
\VUpk{10.2pt}{40}{Die boord.}
\VUsaut{3.4mm}

%%K t03-c1-09-1 p
9. --- Meneer Aliboron (die donkie) is heeltemal tevrede met daardie oggendloop. Hy
snuif met genot die lug in wat deur die geur van die \VUgras{blomme}
\textsuperscript{(20)} balsem is, wat die \VUgras{vrugtebome} \textsuperscript{(18)} van
die naburige \VUgras{boord} bedek. Heeltemal pienk en wit is daardie
\VUgras{perskebome} \textsuperscript{(18)} en \VUgras{appelkoosbome}
\textsuperscript{(19)}, en hulle laat 'n goeie oes hoop.

%%K t03-c1-10-1 p
10. --- Intussen gaan die klein \VUgras{bye} \textsuperscript{(22)} van die
\VUgras{byekorf} \textsuperscript{(21)} daarheen om gonsend hulle soet \VUgras{heuning}
te gaan buit.

%%K t03-c1-11-1 p
11. --- Maar wat gebeur? Hier kom die kinders van die dorp aangehardloop en jaag die
verskrikte \VUgras{ganse} \textsuperscript{(25)} weg. Dit is die uitkoms van die stoet
uit die kerk na die \VUgras{doop} van die notaris se seun.

%%K t03-c1-12-1 p
12. --- Die klein Jakobus, in sy \VUgras{wieg} \textsuperscript{(26)}, word wakker van
die vrolike gelui van die klokke, en sy sussie Magdalena, wat ook die doopstoet wil
sien, vra haar moeder om te kom haar hare kam en haar klere aantrek op die drumpel van
die huis.

%%K t03-c1-13-1 p
13. --- Die moeder stem in. Sy sit haar \VUgras{kopdoek} \textsuperscript{(28)}, haar
\VUgras{lyfie} \textsuperscript{(29)} en haar \VUgras{voorskoot} \textsuperscript{(30)}
aan, en sy kom sit op 'n stoeltjie voor die deur. Magdalena het net haar
\VUgras{onderrok} \textsuperscript{(31)} en haar \VUgras{borsrokkie}
\textsuperscript{(32)}.

%%K t03-c1-14-1 p
14. --- Hoe gelukkig is sy! Sy vergeet haar geliefde blomme, haar \VUgras{angeliere}
\textsuperscript{(34)}, waarop die \VUgras{skoenlappers} \textsuperscript{(35)} gaan
sit, haar \VUgras{tulpe} \textsuperscript{(36)}, haar \VUgras{lelietjies}
\textsuperscript{(37)}, in \VUgras{potte} \textsuperscript{(38)}, op die \VUgras{muur}
\textsuperscript{(33)}, en haar \VUgras{rose} \textsuperscript{(39)} wat so 'n mooi
heining vorm. Sy beskou die ryk klere van die dames van die gevolg.

%%K t03-c1-15-1 p
15. --- Die kinders van die dorp let nie veel op die klere nie. Hulle storm en stamp
mekaar om \VUgras{lekkers} \textsuperscript{(55)} of \VUgras{geldstukke}
\textsuperscript{(52)} te kry. Een van hulle huil \textsuperscript{(40)}.

%%K t03-c1-16-1 p
16. --- 'n Ander een het op die poot van die \VUgras{hond} \textsuperscript{(42)}
getrap, wat kermend wegvlug en die \VUgras{hen} \textsuperscript{(43)} en haar
\VUgras{kuikentjies} \textsuperscript{(44)} wegjaag.

%%K t03-tit-6 sub
\VUsaut{5.0mm}
\VUpk{10.2pt}{40}{Die doopstoet.}
\VUsaut{4.0mm}

%%K t03-c1-17-1 p
17. --- Voor die gevolg loop die \VUgras{min} \textsuperscript{(45)}. Sy het 'n mooi
\VUgras{kappie} \textsuperscript{(46)} met lang rooi linte. Sy dra die
\VUgras{pasgeborene} \textsuperscript{(47)} heeltemal in \VUgras{kant}
\textsuperscript{(48)} geklee.

%%K t03-c1-18-1 p
18. --- Agter die min kom die \VUgras{peetoom} \textsuperscript{(49)} en die
\VUgras{peettante} \textsuperscript{(53)}. Hulle deel die lekkers en geldstukke uit. Die
peetoom het \VUgras{handskoene} \textsuperscript{(51)} en 'n \VUgras{keil}
\textsuperscript{(50)}. Die peettante hou 'n mooi \VUgras{ruiker} \textsuperscript{(54)}
in haar hand.

%%K t03-c1-19-1 p
19. --- Agter hulle sien ons die \VUgras{oom} \textsuperscript{(56)} en die
\VUgras{tante} \textsuperscript{(58)} van die kind. Die tante het 'n elegante
\VUgras{hoed} \textsuperscript{(59)}, die oom 'n swart \VUgras{geklede jas}
\textsuperscript{(57)} en 'n wit onderbaadjie. Hier is dan die \VUgras{neef}
\textsuperscript{(60)} en die \VUgras{niggie} \textsuperscript{(61)}. Hierdie een hou
die \VUgras{sussie} \textsuperscript{(62)} van die pasgeborene, die klein Berta, aan die
hand.

%%K t03-c1-20-1 p
20. --- Die ander \VUgras{boetie} \textsuperscript{(63)} van Berta loop agteraan. Hy gee
sy hand aan 'n \VUgras{familielid} \textsuperscript{(64)}. Die \VUgras{vriende}
\textsuperscript{(65)} van die gesin kom daarna.

%%K t03-tit-7 sub
\VUsaut{4.6mm}
\VUpk{10.2pt}{40}{Die toeskouers.}
\VUsaut{3.4mm}

%%K t03-c1-21-1 p
21. --- Naby die stoet, hier is 'n \VUgras{bedelaar} \textsuperscript{(66)} wat op sy
\VUgras{kruk} \textsuperscript{(67)} steun. Hy vra 'n aalmoes (hy bedel).

%%K t03-c1-22-1 p
22. --- Aan die ander kant is 'n klein seuntjie wat baie \VUgras{beleefd}
\textsuperscript{(68)} is. Hy groet en sê « \VUgras{goeiedag} » vir die mense wat hy
ken.

%%K t03-c1-23-1 p
23. --- Die dorpenaars het uit hulle huise gekom om die doopstoet te sien. Ons sien 'n
\VUgras{boer} \textsuperscript{(69)} met sy \VUgras{katoenmus} en langsaan 'n
\VUgras{boerin} \textsuperscript{(70)}. Dan 'n \VUgras{ou vrou} \textsuperscript{(71)}
wat op haar \VUgras{kierie} \textsuperscript{(72)} steun. Laastens die
\VUgras{meulenaar} \textsuperscript{(73)} met sy groot \VUgras{vilthoed}
\textsuperscript{(74)} en 'n jong boerin met \VUgras{klompe}
\textsuperscript{(75)} aan haar voete, wat haar kindjie leer loop.

%%K t03-c1-24-1 p
24. --- Daardie toneel is baie vermaaklik.
\VUsaut{4.0mm}

%%K t03-tit-8 sub
\VUcentre{12.0pt}{0}{\textit{Tweede tafereel.}}
\VUsaut{3.0mm}

%%K t03-tit-9 sub
\VUpk{10.2pt}{40}{Openbare fees.}
\VUsaut{2.4mm}

%%K t03-tit-10 sub
\VUpk{10.2pt}{40}{Waterspele en roeiwedvaarte.}
\VUsaut{3.0mm}

%%K t03-c2-01-1 p
1. --- Die \VUgras{somer} het gekom. Die warm son van die maand \VUgras{Junie} (Julie of
\VUgras{Augustus}) spoor die jongmense aan om te bad en te roei.

%%K t03-c2-02-1 p
2. --- Op die regteroewer van die rivier trek die roeiers hulle uit om aan die
waterspele en roeiwedvaarte deel te neem.

%%K t03-c2-03-1 p
3. --- Een van hulle trek sy \VUgras{skoene} \textsuperscript{(1)} uit; hy maak hulle
los (hy knoop hulle oop). Sy twee maats is reeds gereed. Van nou af het hulle net hulle
\VUgras{gebreide hemp} \textsuperscript{(2)} en \VUgras{swembroek}
\textsuperscript{(3)}. Hulle gesels terwyl hulle op hulle vriende wag. Hulle praat oor
die kermis en oor die fees wat op die ander oewer plaasvind.

%%K t03-c2-04-1 p
4. --- Op die grond, naby hulle, lê die \VUgras{klere} van hulle maats: 'n paar
\VUgras{stewels} \textsuperscript{(4)}, \VUgras{skoene} \textsuperscript{(5)},
\VUgras{kouse} \textsuperscript{(6)}, \VUgras{vilt-} \textsuperscript{(7)} en
\VUgras{strooihoede} \textsuperscript{(8)} en 'n \VUgras{das} \textsuperscript{(9)}.

%%K t03-c2-05-1 p
5. --- Een van die jongmanne het sy \VUgras{baadjie} \textsuperscript{(10)} sorgvuldig
gevou. Hy sit dit op 'n handdoek, waarin sy buurman gou albei \VUgras{klere} sal
toemaak.

%%K t03-c2-06-1 p
6. --- Die sesde seun is laat. Hy het nog sy \VUgras{pet} \textsuperscript{(11)} op sy
kop. Hy het nie sy \VUgras{hemp} \textsuperscript{(12)}, of sy \VUgras{boordjie}
\textsuperscript{(13)}, of sy \VUgras{mansjette} \textsuperscript{(14)} uitgetrek nie. Hy
hou sy \VUgras{onderbaadjie} \textsuperscript{(17)} in die hand en het nog sy
\VUgras{broek} \textsuperscript{(16)} en \VUgras{kruisbande} \textsuperscript{(15)}. Hy
sal beslis nie gereed wees vir die volgende wedloop nie.

%%K t03-c2-07-1 p
7. --- Naby die jongmanne lei 'n paadjie, waarlangs baie \VUgras{dissels}
\textsuperscript{(18)} staan, na die oewer van die rivier, waar vader Jakobus tussen die
\VUgras{riete} \textsuperscript{(19)} die \VUgras{vuurwerk} \textsuperscript{(20)}
voorberei wat hulle die aand sal afskiet.

%%K t03-tit-11 sub
\VUsaut{4.4mm}
\VUpk{10.2pt}{40}{Die wedloop.}
\VUsaut{3.4mm}

%%K t03-c2-08-1 p
8. --- Op die rivier vaar 'n paar dames, wat hulle met hulle sambrele beskut, in 'n
\VUgras{boot} \textsuperscript{(21)}. Hulle volg die wedloop wat baie interessant is. In
elke \VUgras{skuit} \textsuperscript{(22)} roei vier \VUgras{roeiers}
\textsuperscript{(24)} met al hulle kragte, terwyl die \VUgras{stuurman}
\textsuperscript{(26)} die \VUgras{roer} \textsuperscript{(25)} hou en die brose roeiboot
stuur. Die \VUgras{roeispane} \textsuperscript{(23)} slaan die water op maat en die ligte
bote vaar met duiselingwekkende snelheid.

%%K t03-c2-09-1 p
9. --- 'n Paar jongmanne stry, naby die pilaar van die brug, om die prys van die
\VUgras{seilboom} \textsuperscript{(28)}. Een van hulle loop versigtig op die buigsame en
gladde boom om die klein \VUgras{vlaggie} \textsuperscript{(29)} te bereik wat aan die
punt vasgemaak is. Sy ongelukkige \VUgras{mededinger} \textsuperscript{(30)} het in die
water geval. Hy swem om aan wal te kom.

%%K t03-c2-10-1 p
10. --- Ouer jongmanne \VUgras{hou 'n boottoernooi} \textsuperscript{(32)} naby die
\VUgras{kaai} \textsuperscript{(33)}. Al hierdie spele word bygewoon deur 'n talryke
\VUgras{publiek} \textsuperscript{(31)} wat teen die \VUgras{reling} van die
\VUgras{brug} \textsuperscript{(27)} leun en op die \VUgras{oewer} saamgedrom is. 'n Paar
toeskouers draai hulle egter om, om die spel van die \VUgras{prysmas}
\textsuperscript{(45)} te volg of om die \VUgras{burgemeestersgevolg} te bewonder wat vir
die \VUgras{uitdeling} van die pryse kom.

%%K t03-c2-11-1 p
11. --- Eers is hier 'n paar \VUgras{straatseuns} \textsuperscript{(35)} wat voor die
\VUgras{musikante} \textsuperscript{(36)} loop; dan kom die \VUgras{burgemeester}
\textsuperscript{(37)}, die adjunkte en die \VUgras{gemeenteraad} van die dorpie.
Laastens loop die \VUgras{brandweermanne} \textsuperscript{(38)}.

%%K t03-tit-12 sub
\VUsaut{5.0mm}
\VUpk{10.2pt}{40}{Die kermis.}
\VUsaut{3.6mm}

%%K t03-c2-12-1 p
12. --- Baie mense is op die \VUgras{kermisplein} \textsuperscript{(39)}. 'n Mens kom
die verskillende \VUgras{tente} sien: die \VUgras{houtperde} \textsuperscript{(40)}, die
\VUgras{sirkus} \textsuperscript{(41)} waar die vertoning reeds begin het; die
\VUgras{openbare bal} \textsuperscript{(42)}, waar die jongmanne en die jongmeisies van
die stad die plesier van die \VUgras{dans} geniet.

%%K t03-c2-13-1 p
13. --- Agterin sien ons die \VUgras{arena} \textsuperscript{(44)} waar die dapper
Spaanse \VUgras{stiervegter} teen die woedende \VUgras{bul} \textsuperscript{(45)} veg,
dan die \VUgras{glybaan op spore} \textsuperscript{(49)}, die \VUgras{skietbaan}
\textsuperscript{(46)} waar 'n mens jou oefen om \VUgras{vuurwapens} te gebruik en na die
\VUgras{teiken} te skiet.

%%K t03-c2-14-1 p
14. --- Naby, hier is die \VUgras{kermisteater} \textsuperscript{(47)} waarin die
\VUgras{blyspel} en die \VUgras{treurspel} gespeel word. Baie naby is die tent van die
\VUgras{reus} \textsuperscript{(48)} en van die \VUgras{dwerg}. Die \VUgras{lengte} van
die eerste is twee meter en vyftien sentimeter; die tweede is skaars so groot soos 'n
kind.

%%K t03-c2-15-1 p
15. --- Laastens, hier is die groot \VUgras{dieretuin} \textsuperscript{(50)} met sy
\VUgras{wilde diere}, sy \VUgras{tier} \textsuperscript{(51)} met magtige
\VUgras{kloue}, sy \VUgras{leeu} \textsuperscript{(52)} met 'n digte \VUgras{maanhare},
sy twee \VUgras{kameelperde} \textsuperscript{(53)} en sy \VUgras{olifant}
\textsuperscript{(54)} wat sy \VUgras{slurp} so behendig gebruik.

%%K t03-c2-16-1 p
16. --- Die \VUgras{temmer} \textsuperscript{(55)} wat daardie diere afrig, staan by die
deur van die tent.
\VUsaut{6.0mm}
\VUfilet{22mm}
